\section{Motivation for the Research} \label{section:motivation}

Interfaces rarely form under ideal conditions. Instead, they emerge through complex, often non-equilibrium processes such as high-temperature annealing, epitaxial growth, and chemically driven surface reconstructions. These dynamics frequently trap systems in metastable configurations or generate local disorder, resulting in atomic arrangements that deviate significantly from those predicted by ideal lattice-matching rules. Such complexity challenges efforts to model interface behaviour using conventional approaches alone.

While Density Functional Theory (DFT) provides a reliable route to total energy prediction and structural relaxation, its unfavourable scaling with system size ($\mathcal{O}(N^3)$ in electron count) severely limits its use in high-throughput workflows. In particular, DFT becomes prohibitive when applied to interface supercells exceeding a few hundred atoms, as required to capture long-range strain, multi-layer reconstruction, or compositional variation. Moreover, standard DFT approximations, such as semi-local exchange-correlation functionals, may fail to capture interfacial band alignments, polarisation effects, or weak van der Waals interactions critical to layered systems.
% need to discuss how MLPs provide solutions to the problems of failure to capture interfacial band alignments, polarisation effects, or weak van der Waals interactions critical to layered systems. Explain in the paragraph below...
% MLP basically solves this problem of accurate predictions of interfacial band alignments, polarisation effects, or weak van der Waals interactions critical to layered systems by being able to take into account long range interactions, however given that MLPs are trained on DFT data, they are only as good as the DFT data they are trained on. This means that if the DFT data is not accurate, the MLP will not be accurate either. This is a limitation of MLPs. MLPs can provide a more efficient way to sample the configurational space.
Machine-learned interatomic potentials (MLPs) offer a promising route to overcoming these limitations. By learning to emulate DFT-level energies and forces, MLPs like MACE provide a means to conduct large-scale structure relaxations and energy evaluations at dramatically reduced cost. This enables exploration of broader configurational spaces, including low-symmetry systems and extended defects, that would otherwise remain intractable. However, this advantage is only realised if the models generalise well across diverse interfacial environments.

The core motivation of this work, therefore, is to assess whether MLPs can replicate the DFT-derived ordering of interface favourability, defined here in terms of relaxed formation energy, across a wide range of semiconductor heterostructures. The benchmarking in Chapter~\ref{chapter:computational_investigation} demonstrates that MLPs can often preserve the relative rankings of DFT energies, especially in chemically similar systems. However, notable discrepancies arise in cases such as pure-Si interfaces, where model performance degrades, prompting the need for refinement or retraining. These results motivate the delta-learning strategies explored in Chapter~\ref{chapter:next_steps}, which aim to correct systematic MLP-DFT residuals through supervised learning on interface-specific error distributions.
% not the goal of the PhD goal of what i have been looking at for a while, but the goal of the PhD is create a Machine Learning Interface (MLI) model that can predict the interface energy based on the top and bottom slabs, without having to construct the full interface. This is a long term goal of the PhD, but it is not the goal of this report. I am specifically looking at rapid processing to fast find the best interfaces to train and evaluate the MLPs on, and then use the MLPs to predict the interface energy. This is a long term goal of the PhD, but it is not the goal of this report. I am specifically looking at rapid processing to fast find the best interfaces which can then be processed in DFT ISIF3 a field accepted means on structural evaluation.
In addition to correcting prediction accuracy, this research investigates whether interface favourability can be inferred directly from structural or surface-level features. Chapter~\ref{chapter:next_steps} proposes the use of learned models that map surface properties to interfacial energies, potentially bypassing the need for full supercell construction. This complements the structure-generation pipelines, based on ARTEMIS, that systematically enumerate viable interface candidates across various stacking shifts and Miller plane pairings.

Ultimately, this project is motivated by the need for more efficient and predictive tools to accelerate the design of functional heterostructures. By embedding MLPs within automated screening workflows and augmenting them with delta-learning corrections and structure-based heuristics, it becomes possible to bridge the gap between physical realism and computational tractability. In doing so, this research contributes to the broader aim of enabling scalable, data-driven exploration of complex interfacial systems; particularly in regimes where conventional symmetry-based rules or bulk-derived heuristics break down.